\documentclass[11pt,a4paper]{article}
\usepackage[margin=25mm,headheight=14pt]{geometry}
\usepackage{fontspec}
\setmainfont{Noto Serif}
\setsansfont{Noto Sans}
\setmonofont{DejaVu Sans Mono}[Scale=0.82]
\usepackage{amsmath,amssymb,amsthm,mathtools}
\usepackage{booktabs,longtable,tabularx,array}
\usepackage{enumitem}
\usepackage{fvextra}
\usepackage{xcolor}
\usepackage{microtype}
\usepackage[unicode,hidelinks]{hyperref}
\usepackage{bookmark}
\usepackage{xurl}
\usepackage{titlesec}
\titleformat{\section}{\Large\sffamily\bfseries}{\thesection}{0.7em}{}
\titleformat{\subsection}{\large\sffamily\bfseries}{\thesubsection}{0.7em}{}
\titleformat{\subsubsection}{\normalsize\sffamily\bfseries}{\thesubsubsection}{0.7em}{}
\setlength{\parskip}{4.5pt}
\setlength{\parindent}{0pt}
\setlength{\emergencystretch}{2em}
\setlist{itemsep=2pt,topsep=4pt}
\DefineVerbatimEnvironment{Code}{Verbatim}{breaklines=true,fontsize=\small,baselinestretch=1.03}
\newcommand{\code}[1]{\texttt{#1}}
\newcommand{\R}{\mathbb R}
\newcommand{\Q}{\mathbb Q}
\newcommand{\N}{\mathbb N}
\newcommand{\Z}{\mathbb Z}
\newcommand{\FV}{\operatorname{FV}}
\newcommand{\eval}{\operatorname{eval}}
\newcommand{\app}{\operatorname{app}}
\newcommand{\rev}{\operatorname{rev}}
\newcommand{\racc}{\operatorname{revAcc}}
\newcommand{\supp}{\operatorname{supp}}
\newtheorem{proposition}{Proposition}[section]
\newtheorem{definition}[proposition]{Definition}
\newtheorem{invariant}[proposition]{Invariant}
\newtheorem{example}[proposition]{Example}
\hypersetup{pdftitle={Forge: A Certificate-Driven Proof Planner Beyond Lean's grind},pdfsubject={Concrete tactic architecture, algorithms, independent prototypes, and evaluation},pdfauthor={}}
\begin{document}
\begin{center}
{\LARGE\sffamily\bfseries Forge: A Certificate-Driven Proof Planner\\[5pt] Beyond Lean's \texttt{grind}}\par
\end{center}
\vspace{4mm}
\begin{abstract}
A substantially more capable general-purpose Lean tactic should not replace
\code{grind}'s effective local reasoning with a collection of disconnected
solvers. It should discover the proof structure that makes local reasoning
succeed: the right induction motive, an auxiliary lemma, an existential witness,
a missing theorem instance, a nonlinear inequality certificate, or a necessary
representation bridge. This article specifies \emph{Forge}, a proposed
certificate-driven proof planner that retains \code{grind} as a central engine
and connects these capabilities through a typed, scope-aware graph of facts and
obligations.

The proposal includes concrete search algorithms, certificate formats, Lean
integration boundaries, a dependency and trust policy, worked mathematical
examples, and an implementation roadmap. Five mechanisms are prototyped in
Python: demand-directed Horn reasoning, structural induction with accumulator
lemma synthesis, sparse rational nonnegativity certificates, Bernstein box
certificates, and polynomial witness synthesis. The accompanying suite passes
69 tests; 13 persisted certificate bundles replay using only Python's standard
library. These results validate small algorithmic fragments, not a completed
Lean tactic or a measured advantage over \code{grind}. The Lean reference
scripts are supplied but were not compiled because a Lean runtime was unavailable.
\end{abstract}

\section*{Executive assessment}
\addcontentsline{toc}{section}{Executive assessment}
\textbf{Recommendation.} Implement Forge as a proof planner with a shared
obligation graph, not as a new monolithic decision procedure. Its first major
capability gain should be \emph{automatic structural induction plus checked
lemma discovery}. Its second should be \emph{witness and theory-certificate
search driven by the needs of the current proof}. Reuse existing Lean arithmetic,
superposition, simplification, and sum-of-squares infrastructure wherever
possible.

The critical distinction is between a proposed fact and a proved fact. Search
engines may be heuristic, numerical, learned, or external. They may request
additional assumptions, but they must not silently receive them. Only a proof
of the exact requested statement, under explicitly recorded hypotheses, may
enter the fact store. Completion means a closed proof of the original Lean
goal with an approved axiom dependency set.

\textbf{Status of this artifact.} The architecture and Lean-facing interfaces
are a design. The Python algorithms, tests, benchmark runner, saved certificates,
and independent replay command are implemented and executed. No end-to-end
\code{forge} tactic, Lean-kernel port of the Python checkers, or comparative Lean
benchmark is claimed. The dated source audit uses Lean 4.34.0 and sources observed
on September 14, 2026; mutable third-party sources are explicitly distinguished
from that pinned baseline.

\clearpage
{\small\setlength{\parskip}{0pt}\tableofcontents}
\clearpage

\section{The actual baseline: what must not be reinvented}
\label{sec:baseline}

\subsection{A versioned comparison, not a caricature}
The starting reference is the Lean language manual's chapter on \code{grind}.
It describes a proof-producing, SMT-inspired combination of congruence closure,
propagation, quantified instantiation, case splitting, and theory reasoning.
Its shared-state organization is already a major part of its strength.
Adding an e-graph, arithmetic, or E-matching is therefore not, by itself, a
meaningful advance over the baseline.\footnote{Lean project, \emph{The grind
tactic}, language reference \cite{grind}.}

The release page resolved to Lean \code{v4.34.0} during this audit, with a
September 14, 2026 release date. The source-level baseline in this article is
that tag, rather than an unspecified future meaning of \code{latest}.
The executable was not available locally, so statements about its implementation
are source observations, not newly run success or failure experiments.\footnote{Lean
project, release \code{v4.34.0} \cite{release}.}

The pinned configuration enables more than the minimal description suggests:
extensionality, function extensionality and function-valued congruence,
associative--commutative reasoning, order reasoning, model-based theory
combination, integer and linear arithmetic, and commutative-ring reasoning.
There is also an optional library-suggestion path. Selected resource defaults
are five E-matching rounds, generation eight, 1,000 instances, and nine splits.
These are budget parameters, not semantic limits of Lean.\footnote{Lean 4.34.0,
\code{Init/Grind/Config.lean} \cite{config}.}

The ring engine has tracked derivations for polynomial consequences, including
superposition and simplification operations. Its proof objects retain algebraic
side conditions associated with cancellation and characteristic. It is
incorrect to describe it merely as a syntactic ring normalizer and then claim
that adding Gr\"obner-style reasoning would be a new capability.\footnote{Lean
4.34.0, commutative-ring constraint and derivation types \cite{commring}.}

The integer-arithmetic reference describes a linear fragment with inequalities,
equalities, disequalities, and divisibility, including natural-number
translation. Nonlinear expressions may be abstracted as atoms in the linear
solver. A satisfying assignment to that abstraction is not necessarily a model
of the original nonlinear formulas.\footnote{Lean language reference,
\emph{Linear Integer Arithmetic} \cite{lia}.}

\subsection{Where a broader tactic can make a real difference}
The opportunity is principally \emph{proof organization}. A local solver may
succeed once a human has chosen an induction, strengthened the induction
hypothesis, supplied a useful identity, constructed a witness, selected the
right theorem, or exposed the right representation. Forge should automate those
choices and feed their checked consequences back into local reasoning.

This is a capability target, not a claim that every worked example below defeats
current \code{grind}. Imported theorems, attributes, preprocessing, and tuning
can change which tactic solves a particular goal. An example involving list
reversal can become trivial when its correctness theorem is already available.
A fair test must specify both the theorem statement and the environment in
which it is presented.

\begin{longtable}{p{0.25\linewidth}p{0.32\linewidth}p{0.35\linewidth}}
\toprule
Capability & Baseline to retain & Forge's proposed addition \\
\midrule
\endhead
Local equality and arithmetic & Existing \code{grind} theory engines & Request and reuse missing structural or nonlinear consequences. \\
Quantified reasoning & Existing E-matching and optional suggestions & Backward demand propagation, controlled new-term generation, and premise subgoals. \\
Recursive functions & Defining equations and local reasoning & Choose induction motives; generalize changing parameters; synthesize and prove auxiliary lemmas. \\
Existential conclusions & Any applicable existing logical rules & Search typed witness grammars and certify their universally quantified obligations. \\
Nonlinear inequalities & Existing arithmetic tactics and external SOS library & Reuse certified engines, coordinate guards, prioritize sparse certificates, and add box reasoning. \\
Trust and diagnostics & Lean proofs and solver traces & Exact-goal replay, explicit unresolved obligations, dependency audit, and reproducible proof artifacts. \\
\bottomrule
\end{longtable}

\subsection{The meaning of ``more powerful''}
Three different claims must be separated. \emph{Broader coverage} means solving
additional families of goals with comparable human input. \emph{Better search}
means solving more instances under a fixed resource budget. \emph{Faster replay}
means a smaller or cheaper proof after search has finished. One can improve
coverage while worsening latency on easy goals.

Forge should expose a baseline-preserving configuration: run the original
\code{grind} configuration before escalating, and retain its successful proof.
With that original budget \emph{plus additional resources}, this preserves
its successes. Under an unchanged total wall-clock budget, no such dominance
follows: orchestration overhead and budget allocation can lose cases. A
non-starving unbounded schedule that eventually reruns every baseline attempt
preserves the corresponding eventual successes, but that observation is not a
practical performance theorem.

\section{Existing libraries and the boundary of the contribution}
\label{sec:libraries}

\subsection{Proof-search and external-prover components}
Aesop already provides extensible proof search organized around goals and rules.
Its AND/OR-style search and distinction between safer and branching rules are
useful precedents and reusable infrastructure. Forge's proposed distinction is
the exchange of \emph{typed demands and certified theory consequences} across
that search, rather than the mere existence of backward search.\footnote{Aesop
project documentation \cite{aesop}.}

Duper is a Lean-integrated proof-producing prover, and a natural candidate for
logical and superposition reasoning. Its fact-selection interface matters:
\code{duper [*]} requests local facts, whereas bare \code{duper} should not be
assumed to consume them automatically. A Forge adapter should pass an explicit,
recorded premise slice.\footnote{Duper repository documentation \cite{duper}.}

Lean-auto provides a translation and automation interface between Lean and
external reasoning. Its README distinguishes reconstruction-capable paths from
standalone SMT or TPTP invocation modes without reconstruction. An external
\code{unsat} result is therefore not a uniform proof interface. Forge must
accept a reconstructed proof or create explicit remaining obligations; it must
not promote a solver answer to a theorem.\footnote{Lean-auto repository and
interface paper \cite{auto,autopaper}.}

Lean-SMT reconstructs cvc5 reasoning in Lean, but its documented interface can
leave proof holes as Lean goals. That is useful partial progress, not completed
verification. The adapter should import those goals as children of the current
AND node and report success only after every child is discharged.\footnote{Lean-SMT
repository \cite{smt}.}

QuerySMT is particularly close prior work: it uses SMT-generated hints together
with Lean reasoning, including \code{grind} and Duper. Combining these systems
is not a new proposal here. Forge's additional design target is a persistent
planner that can generate structural lemmas and witnesses, solve the resulting
obligations, and let those results influence future search.\footnote{Clune,
Barbosa, and Avigad, \emph{Hint-Based SMT Proof Reconstruction} \cite{querysmt}.}

\subsection{Nonlinear arithmetic: reuse the existing SOS engine}
A current \code{leanprover/sos} library already implements an end-to-end
nonlinear-real-arithmetic tactic. Its documented pipeline includes reification,
CSDP search, rational reconstruction, and certificate-based proof construction;
it also provides a witness-oriented replay interface. Consequently, ``add an
SOS tactic to Lean'' is not a novel component of Forge. The production design
should reuse this library after compatibility and trust checks.\footnote{Lean
SOS project, README \cite{sos}.}

The inspected \code{SOS/Engine.lean} exposes \code{SOS.Engine.solve}, a structured
problem/result interface, and deterministic result checking. Its configuration
already includes degree deepening, constraint-product bounds, and optional
refutation powers. Some README wording about fixed relaxation levels lags this
source: the implementation is the stronger basis for planning integration.
The separate \code{SOS.Core} module separates proof-facing data/checking from
native search dependencies.\footnote{Lean SOS source, \code{Engine.lean} and
\code{Core.lean} \cite{sosengine,soscore}.}

There is a real compatibility task, not a ready-made common workspace. The
retrieved SOS toolchain file specifies \code{v4.32.0-rc1}, whereas the inspected
Lean baseline is 4.34.0. That retrieved file may itself be a cached snapshot.
Resolve commits and manifests together before attempting an integration; do
not assume all projects' moving \code{main} branches build with the same Lean
release.\footnote{Retrieved SOS \code{lean-toolchain} \cite{sostoolchain}.}

For sparse polynomial data, \code{Hex.MvPoly} is a concrete candidate substrate.
Its repository describes fixed-arity canonical sparse polynomials with an
explicit monomial order and proved operation/evaluation laws. Forge should
reuse a compatible proof-facing representation rather than repeatedly expand
large expressions into unreduced Lean syntax.\footnote{Hex multivariate
polynomial library \cite{hex}.}

\subsection{Other components and non-goals}
Mathlib's \code{linarith} infrastructure constructs proof terms from arithmetic
certificates; \code{nlinarith} adds useful nonlinear preprocessing. These are
first-line closers, not mechanisms to discard.\footnote{Mathlib,
\code{Mathlib.Tactic.Linarith.Frontend} \cite{linarith}.}
Its \code{positivity} extensions cover compositional sign facts for arithmetic
operations and other supported expressions, and should be tried before a
costly numerical search.\footnote{Mathlib positivity extensions \cite{positivity}.}

The \code{lean-egg} repository is now marked as no longer developed and directs
users toward \code{grind}. It should not be the foundation of a new maintained
Lean package. The Rust \code{egg} library remains a possible untrusted search
component, but importing its data structures does not solve Lean's binder,
dependent-type, or proof-reconstruction problems.\footnote{The \code{lean-egg}
and \code{egg} repositories \cite{leanegg,egg}.}

Induction and lemma synthesis are established research topics, not inventions
of this article. Directed lemma synthesis for functional-program equivalence
is a relevant source of stronger search heuristics. The present contribution
is a concrete integration design and a small, inspectable executable study,
not a claim to have originated these techniques.\footnote{Sun et al.,
\emph{Proving Functional Program Equivalence via Directed Lemma Synthesis}
\cite{lemmasynthesis}.}

\section{A precise architecture for Forge}
\label{sec:architecture}

\subsection{The planner's state}
A search session has the conceptual state
\[
 \mathcal S=(E,\Gamma,G,\mathcal O,\mathcal F,\mathcal D,\mathcal C,\mathcal Q,B).
\]
Here $E$ is an immutable environment snapshot; $\Gamma$ is a typed local-context
snapshot; $G$ is the original target; $\mathcal O$ is an AND/OR obligation graph;
$\mathcal F$ is a store of proved facts; $\mathcal D$ is a store of demands;
$\mathcal C$ contains normalized terms and theory-specific representations;
$\mathcal Q$ is a collection of work queues; and $B$ is a multidimensional
resource budget. ``Immutable'' refers to a snapshot identity: a branch can
produce a successor snapshot, but another branch does not mutate it in place.

A fact record stores its proposition, proof, source scope, dependency support,
and originating certificate identifier. A demand record stores a desired
proposition or term shape, the consumer requesting it, the local context in
which it is needed, and a bounded search plan. A demand carries \emph{no logical
authority}. Two consumers may share one demand only after their scopes and
instantiations have been reconciled.

An obligation node is either a goal with alternative proof methods, an AND
node requiring several subproofs, or a completed proof. A method such as
applying a theorem creates an AND node for its premises. Choosing one of two
theorems creates an OR choice. Induction creates an AND node for the recursor's
minor premises; an existential witness creates an AND node for its remaining
properties.

\subsection{A small vocabulary of engine messages}
The following is an interface specification, not existing Lean API syntax:
\begin{Code}
EngineMessage :=
  Certified(proposition, proof, scope, dependencies)
| Need(proposition, scope, consumer)
| Refine(goal, subgoals, proofBuilder, savedContext)
| Candidate(certificateData, expectedStatement, obligations)
| Conflict(assumptions, checkedRefutation)
| Unknown(reason, resumableSearchState)
\end{Code}

\code{Candidate} is intentionally separate from \code{Certified}. A numerical
optimizer may return a polynomial decomposition, but its equality, sign
conditions, interpretation map, and connection to the original Lean expression
must be checked before a fact is emitted. An SMT engine may return a useful
implication whose premises become demands. A failed external run emits
\code{Unknown}, not a logical negation.

The proof builder in \code{Refine} is a typed Lean function or replay procedure
that combines proofs of its exact subgoals into a proof of its parent. It is not
an unchecked promise that some later tactic will make the relationship true.
Before a method is admitted to the planner, its parent/child relationship must
be elaborated in its saved context.

\subsection{Scope and dependency discipline}
Let a branch have assumptions $A_1,\ldots,A_k$. If a solver proves $P$ using
$A_2$ and $A_5$, the exported fact has support $\{2,5\}$ together with all
necessary local variables and their type dependencies. It may be reused in a
descendant containing that support. It may be exported upward only as a
corresponding implication, after the proof is abstracted over those assumptions.
It is not valid to copy the bare proposition $P$ into a sibling branch.

Computing support requires the dependency closure of free variables in both
the proof and the types of those variables. This matters in Lean: a term
\code{x : Fin n} depends on \code{n}, and a local proof may mention a value only
through a type. A syntactic set of visible hypothesis names is insufficient.
Universe parameters and local typeclass instances are also part of the context.

\begin{invariant}[Certified fact invariant]
Every fact visible at a node has a checked proof in that node's local context,
or a checked abstraction that is explicitly instantiated there. No candidate,
relaxation model, sampled equation, or unresolved solver hint is a fact.
\end{invariant}

\begin{invariant}[Exact parent invariant]
Every completed obligation contains a proof of its recorded target in its
recorded context. A goal may be normalized only together with a proof-producing
transport back to the previous target.
\end{invariant}

\begin{proposition}[Compositional soundness of the planner design]
Assume the engine adapters produce valid Lean proofs or sound proof builders,
and the two invariants are maintained. A completed root yields a Lean proof of
$G$ under $\Gamma$. Untrusted search choices cannot invalidate this conclusion.
\end{proposition}
\begin{proof}
Proceed in a topological order through the completed proof DAG. A leaf is a
checked fact or a checked direct proof. At an AND node, apply the checked proof
builder to the already justified children. At an OR node, use the one completed
alternative. Scope abstraction and instantiation preserve the stated contexts.
The exact-parent invariant transports the final result to $G$. This is a
specification-level argument: correctness of an eventual implementation still
requires Lean checking of the assembled term and an audit of its dependencies.
\end{proof}

\subsection{Why this is not a sequential tactic wrapper}
Consider a goal that requires a bound on a recursive accumulator. Induction
produces a step obligation. A theorem instance suggests an inequality whose
premise is a polynomial nonnegativity claim. The nonlinear engine proves that
claim, which enables the theorem instance, which closes the induction step.
The resulting inductive lemma then simplifies an existential-witness obligation.
A list of tactics run once in sequence has no explicit representation of this
feedback loop or of the fact that several consumers need the same intermediate
claim.

Forge's planner persists those obligations, schedules their producers, and
reactivates waiting consumers when a proof arrives. A sophisticated existing
tactic could implement such behavior internally; the point is to make the
coordination protocol explicit and reusable across independent engines.

\section{Demand-directed quantified reasoning}
\label{sec:demand}

\subsection{Backward requests coupled to forward saturation}
Suppose the goal requires $P(t)$ and a theorem has the form
\[
 \forall \bar x,\quad H_1(\bar x)\to\cdots\to H_m(\bar x)\to P(f(\bar x)).
\]
Match the conclusion against the requested atom, modulo already proved
equalities. A successful substitution $\theta$ creates demands for
$H_1\theta,\ldots,H_m\theta$. When all are proved, instantiate the theorem and
publish $P(f(\bar x))\theta$. This can construct intermediate terms that a
purely ground-driven process has not yet produced.

The forward side remains essential. New proved equalities can enable previously
blocked matches, and facts not directly anticipated by a goal can expose a
contradiction. Use forward saturation as one queue and backward demands as
another, with bounded widening rather than a permanent relevance filter.
Requests are indexed by predicate/function head, type, and selected normalized
argument features; theorem conclusions use a discrimination-tree index.

For multi-premise rules, order joins by estimated selectivity and share partial
matches. Store which premise instances have already been tried. Newly arriving
facts extend only joins that mention their indexes, in the spirit of semi-naive
evaluation. A failed join is not cached forever if it can be enabled by a later
equality merge.

\begin{Code}
solveDemand(target, context, fuel):
    if a scoped proof of target exists: return it
    if fuel exhausted: return Unknown
    for theorem in conclusionIndex(target):
        for substitution in boundedTypedMatches(theorem, target):
            premises := instantiatePremises(theorem, substitution)
            create shared child demands for premises
            when every child is proved:
                replay theorem application
                publish the checked target proof
    schedule a fair forward-saturation slice
\end{Code}

\subsection{Variables absent from the conclusion}
Backward matching does not determine every theorem parameter. In
$\forall x\,y,\ R(x,y)\to P(x)$, a demand for $P(a)$ leaves $y$ unknown. Forge
must not silently choose an arbitrary inhabitant. It can enumerate appropriately
typed terms from the local context, ask a witness engine, exploit equalities,
or leave the application suspended with an existential-style obligation.
Every selected instantiation is checked by Lean.

The executable Horn prototype deliberately handles only rules whose backward
premises become ground after matching the head. It abstains on the example
above even when a witness happens to exist. This sharply bounded fragment makes
its behavior and certificate checker easier to audit; it is not the full
quantifier algorithm proposed for Lean.

\subsection{A measured branching example}
Use an uninterpreted type, constant $a$, unary functions $f,g$, and rules
\[
 P(a),\qquad P(x)\to P(f(x)),\qquad P(x)\to P(g(x)).
\]
The target is $P(f^d(a))$. The included forward baseline expands a breadth-first
frontier and stops as soon as the target is produced. With the supplied rule
order, it generates exactly $2^d$ distinct facts for $d\ge1$: the preceding
levels contain $2^d-1$ facts, and the first target at the next level adds one.
A backward proof uses $d+1$ fact nodes.

This is a separation of \emph{these two implemented Horn search policies}.
It is not a measurement of \code{grind}, whose goal internalization, congruence
classes, triggering, heuristics, and resource policy are not modeled by this
baseline. The experiment demonstrates why demand information can matter; it
does not establish a Lean speedup.

The logical proof has linear size in $d$, but the Python implementation uses
nested term objects. Hashing and comparing such terms can add depth-dependent
costs. A production hash-consed term DAG is needed before translating the node
count into an approximately linear implementation cost.

\subsection{Proof certificates and termination controls}
A Horn certificate is a topologically ordered DAG. Each node is either an
explicit premise or a named rule instance whose parent indices are smaller
than the node index. The checker independently matches the rule's head and all
premises under one consistent substitution. The final node must be the exact
requested goal.

Cycle detection prevents immediate recursive demand loops. It is not a
completeness argument for arbitrary recursive Horn programs: an active request
may have another finite derivation that a particular bounded policy misses.
The planner should use iterative depth bounds, table successful answers, and
retain resumable alternatives. A timeout or active-cycle rejection must be
reported as absence of a certificate, not as falsity.

\section{Automatic induction and auxiliary-lemma discovery}
\label{sec:induction}

\subsection{Selecting an induction variable and motive}
Induction should be an explicit planner action, not an unconditional tactic
applied to every natural number or list. Analyze unresolved recursive
applications after cheap simplification. For each application, inspect the
defining equations and identify structurally decreasing arguments, parameters
that change in recursive calls, and dependencies of the target on those
arguments.

Generate a bounded collection of plans: ordinary structural induction on a
visible datatype variable; functional induction using the definition's available
induction principle; and, later, well-founded induction with a supplied or
proved decrease relation. Initially prefer one major premise. Trying arbitrary
nested inductions before examining the residual goal is a recipe for a large
search tree.

A useful ranking function rewards recursive occurrences that become constructor
applications and penalizes the number of constructors, dependent premises, and
new generalized variables. Ranking affects only search order. The actual
induction step must be performed by applying a Lean recursor or a checked
induction theorem; Forge must not manufacture an unrestricted ``induction
hypothesis'' as an ordinary axiom.

The motive is the main source of additional capability. If a recursive call
changes an accumulator, an induction hypothesis about one fixed accumulator may
be useless. Generalize that parameter before induction. In a dependent context,
also revert hypotheses whose types depend on a generalized variable, in a
valid dependency order, and reintroduce them in each minor premise. The target
then quantifies over the appropriate parameters and their assumptions.

\subsection{A fully worked accumulator example}
Write list concatenation as $\app$. Define
\begin{align*}
 \app([],a)&=a,&
 \app(x::xs,a)&=x::\app(xs,a),\\
 \rev([])&=[],&
 \rev(x::xs)&=\app(\rev(xs),[x]),\\
 \racc([],a)&=a,&
 \racc(x::xs,a)&=\racc(xs,x::a).
\end{align*}
The original target is
\begin{equation}\label{eq:original}
 \forall xs,\quad \racc(xs,[])=\rev(xs).
\end{equation}
A direct induction yields an induction hypothesis about
$\racc(xs,[])$, but the step contains $\racc(xs,[x])$. The desired
strengthening is
\begin{equation}\label{eq:strengthening}
 \forall xs\,a,\quad \racc(xs,a)=\app(\rev(xs),a).
\end{equation}

The implemented prototype does not simply return the right-hand side of
\eqref{eq:strengthening}. It notices that the second argument of the recursive
\code{revAcc} call changes, replaces the original fixed accumulator by a
variable, and enumerates a small right-hand-side grammar:
\[
 e ::= xs\mid a\mid []\mid \rev(e)\mid \app(e,e).
\]
Terms are enumerated in increasing syntax-tree size. There are 49 concrete
sample environments: every pair of lists of length at most two over a two-element
alphabet. Samples only reject candidates. The first surviving candidate that
receives an induction certificate is the 41st enumerated term,
$\app(\rev(xs),a)$.

The actual generalization recognizer is deliberately specific to the
\code{revAcc} equation. The general motive-analysis algorithm above is a design
for the Lean implementation, not an already implemented generic recursion
analyzer. This distinction prevents a tiny successful example from being
misrepresented as general automatic induction.

\subsection{The resulting proof, not just its tests}
First prove concatenation's right identity and associativity by induction.
These are named, separately checked lemmas in an acyclic theorem bank. The base
case of \eqref{eq:strengthening} is
$\racc([],a)=a=\app([],a)$.
For the step, assume
\[
 \forall a,\quad\racc(xs,a)=\app(\rev(xs),a).
\]
Then
\begin{align*}
 \racc(x::xs,a)
 &=\racc(xs,x::a)\\
 &=\app(\rev(xs),x::a)\\
 &=\app(\app(\rev(xs),[x]),a)\\
 &=\app(\rev(x::xs),a).
\end{align*}
The middle equality uses associativity and the defining equation for appending
a singleton. Specializing the proved lemma at $a=[]$ and using right identity
proves \eqref{eq:original}.

The prototype independently checks the base and step rewrite traces. In the
step, the tail is a \emph{rigid eigenconstant}; only explicitly generalized
parameters remain instantiable in the induction hypothesis. In the base case,
the induction hypothesis is not available at all. These restrictions are
essential. Treating the tail as an ordinary rewrite metavariable would allow
an induction hypothesis to rewrite the original goal at an unjustified larger
argument.

\subsection{From failed residuals to stronger lemma synthesis}
A production synthesizer should use failure information instead of enumerating
all small theorems. Normalize both sides of an unclosed induction step and
identify the unmatched recursive skeletons. Compute typed anti-unifiers of the
original target and those residuals, recording substitutions back to each.
Candidate lemmas should relate the same operations while replacing obstructing
subterms by fresh parameters.

There are three useful candidate sources. First, generalize a changing
non-structural argument, as above. Second, propose a homomorphism/distribution
lemma such as $\rev(\app(xs,ys))=\app(\rev(ys),\rev(xs))$ when nested calls
obstruct normalization. Third, synthesize a stronger invariant containing a
measure, bound, or relation required by the step. Each candidate becomes its
own proof obligation, never a rewrite rule merely because it helped simplify
another candidate.

Use a typed, bounded grammar assembled from operations present in the residual,
relevant constructors, and a small retrieved lemma vocabulary. Eliminate
candidates that fail executable tests when testing is available. Quotient
syntactically different candidates only by \emph{already proved} equalities;
otherwise the quotient can discard the correct expression or conceal an
unproved assumption. After a failed universal proof, retain the candidate as
unproved; a failure need not yield a counterexample.

A stronger future implementation can replace broad enumeration with the
induction-friendly transformations studied in directed lemma synthesis. That
is an algorithmic development path, not a claim that the included enumerator
implements that research system.

\subsection{Preventing circular discoveries}
Each candidate lemma has a dependency set. It may use previous accepted lemmas
and its own legally scoped induction hypotheses. It may not use itself as a
global theorem, the original unproved target, or another candidate that depends
back on it. The simplest initial policy is an append-only theorem bank ordered
by accepted proofs. More elaborate mutually inductive schemes require an
explicit joint induction principle, not a cycle in a generic dependency graph.

The prototype's bank checks five list lemmas: concatenation right identity,
concatenation associativity, the discovered accumulator specification,
reversal over concatenation, and reversal involutivity. Only the accumulator
specification's right-hand side is synthesized. The other four are supplied
statements whose induction certificates are produced and replayed. All are
about a small freely generated list fragment, not arbitrary dependent Lean
objects.

\section{Witness synthesis connected to proof obligations}
\label{sec:witness}

\subsection{A witness is a program with a proof obligation}
For a goal $\exists y:\tau,\ P(y)$, create a typed grammar for $y$ from the
local context and relevant constructors. The alternatives are not just constant
terms: after universal variables have been introduced, a witness may depend on
them. Instantiating a candidate $w$ produces the exact child goal $P(w)$.
Success requires a proof of that child, followed by the ordinary existential
introduction rule.

For a goal $\forall x,\exists y,\ P(x,y)$, sampling individual values of $x$
and finding unrelated values of $y$ does not construct a function or prove the
theorem. The grammar must represent a well-typed expression $w(x)$, or the
planner must produce another legitimate Lean proof of the existential statement.
A synthesis oracle may help choose $w$ but cannot justify $P(x,w(x))$ by testing.

\subsection{An executed polynomial-witness example}
Consider
\[
 \forall x\in\R,\quad \exists y\in\R,\quad x\le y\ \wedge\ -x\le y.
\]
The prototype enumerates $w(x)=a+bx+cx^2$ with
$a,b,c\in\{-2,-1,0,1,2\}$, ordered first by coefficient $\ell_1$ size and then
lexicographically. Five values $x=-2,-1,0,1,2$ are used only as a rejection
filter. The 23rd candidate is $w(x)=x^2+1$; three candidates reach universal
certificate search before acceptance.

The certificate engine separately proves $w-x\ge0$ and $w+x\ge0$ over all real
inputs. A compact mathematical presentation is
\[
 x^2+1-x=(x-\tfrac12)^2+\tfrac34,\qquad
 x^2+1+x=(x+\tfrac12)^2+\tfrac34.
\]
These are convenient explanatory decompositions. The numerical LP run actually
returned different rational weighted-square decompositions, retained verbatim
as data in the certificate file; both replay exactly. The artifact does not
pretend that a hand-simplified certificate is the optimizer's original output.

The example illustrates cross-component behavior: witness enumeration requests
two universal nonlinear facts, and a theory certificate closes them. A
production implementation should also try simpler witnesses such as an
available absolute-value or maximum expression when the local library makes
those cheap. Restricting the prototype to polynomials is an experiment design,
not a recommended restriction on all Forge witnesses.

\subsection{Guards and dependent witnesses}
A candidate involving division, an array lookup, or a dependent constructor
may require additional evidence. That evidence becomes an explicit child goal.
For example, constructing a value of \code{Fin n} requires a value and a bound;
constructing a term with a nonzero-denominator simplification requires the
appropriate nonzeroness statement. Shared guard demands should be solved once
and reused with their proofs.

The planner should distinguish candidate elaboration failure from mathematical
refutation. It should also preserve dependencies between witness metavariables:
a later witness may depend on an earlier one. Eagerly closing a goal using a
metavariable that another branch later instantiates is unsafe as a search-state
operation even when Lean would ultimately reject the resulting term.

\section{Exact nonlinear reasoning and certificate search}
\label{sec:nonlinear}

\subsection{The certificate language}
Normalize a polynomial target over an explicitly selected ordered field.
Let $g_j\ge0$ be proved inequality guards and $h_k=0$ proved equality guards.
A general certificate can have the form
\begin{equation}\label{eq:cone}
 p = \sum_{J}\sum_i c_{J,i}q_{J,i}^2\prod_{j\in J}g_j
       +\sum_k u_k h_k,
 \qquad c_{J,i}\in\Q_{\ge0}.
\end{equation}
The finite index $J$ can be a set or a multiset, depending on the search policy.
The polynomials $q_{J,i}$ and $u_k$ have rational coefficients. The equality is
checked exactly, not by floating-point residuals.

\begin{proposition}[Conditional cone certificate]
If \eqref{eq:cone} holds as a polynomial identity and the stated guards hold at
an input, then $p\ge0$ at that input.
\end{proposition}
\begin{proof}
Each square is nonnegative, each coefficient is nonnegative, and each guard
product is nonnegative. Their sum is nonnegative. Every equality-ideal term
$u_kh_k$ evaluates to zero. Evaluating the polynomial identity gives the claim.
\end{proof}

This simple argument is the core reason an unreliable optimizer can be useful.
It searches for data; a small checker verifies a mathematical identity and a
finite list of sign conditions. The typed interpretation from rational
polynomials to the original Lean field is part of the proof, not an implicit
casting convention.

The Python certificate checker implements the nonnegative-cone part, without
equality-ideal multipliers. It accepts arbitrary sparse polynomials as square
arguments, but its search grammar is much more restricted. A successful check
is conditional on the supplied guards; the checker does not prove those guards
from a Lean context.

\subsection{The implemented sparse LP search}
Let $D=\deg p$. For each guard product $g_J$ of bounded length, enumerate
monomials of degree at most
\[
 d_J=\left\lfloor\frac{D-\deg g_J}{2}\right\rfloor.
\]
Create a dictionary containing each monomial $m$ and each binomial
$m+cn$ for distinct monomials and
$c\in\{-2,-1,-\tfrac12,\tfrac12,1,2\}$. The corresponding candidate columns
are $q^2g_J$. Flatten their rational coefficients into a linear system
\[
 Aw=b,\qquad w\ge0.
\]
The prototype asks SciPy's HiGHS LP interface to minimize $\sum_iw_i$ subject
to these constraints.\footnote{SciPy documentation,
\code{linprog(method="highs")} \cite{highs}.}

The returned floating-point weights are rationalized with a denominator cap
of $10^6$. Tiny weights may be dropped during \emph{search reconstruction};
this never weakens the final check. The reconstructed identity must match the
target exactly using \code{fractions.Fraction}. A residual of $10^{-12}$ is
not a proof. If exact replay fails, the result is \code{Unknown} even when the
numerical solver reports success.

With $M=\binom{n+d}{d}$ monomials, the unguarded dictionary has at most
$M+6\binom M2$ entries before deduplication. Guard products increase this count
combinatorially. Sparse coefficient storage, incremental degree bounds,
coefficient-size limits, and proof-size limits are therefore necessary. The
LP prototype is intentionally small; it is not a replacement for the existing
Lean SOS library's more general SDP machinery.

\subsection{An executed quartic certificate}
For
\[
 p=x^4+y^4+z^4-x^2y^2-y^2z^2-z^2x^2,
\]
the search constructs 280 candidate columns and 35 coefficient equations.
It finds a three-term certificate:
\begin{equation}\label{eq:quartic}
 p=\tfrac12(x^2-y^2)^2+\tfrac12(y^2-z^2)^2+\tfrac12(z^2-x^2)^2.
\end{equation}
The saved certificate is checked by exact coefficient comparison. This is a
useful compact replay target, but it is not evidence that other existing Lean
tactics cannot solve the theorem. In fact, supplying the three square facts
makes the final arithmetic proof straightforward; the automation question is
how to find and use those facts without human intervention.

Other executed examples include $xy\ge0$ from $x,y\ge0$, $x^3\ge0$ from
$x\ge0$, $x-x^2\ge0$ from $x,1-x\ge0$, and
\[
 x^3+y^3-xy(x+y)=(x+y)(x-y)^2\ge0
 \quad\text{when }x,y\ge0.
\]
These exercise conditional guard products and squares, not just unguarded
polynomial normalization.

\subsection{A deliberate incompleteness control}
The true polynomial $(x+y+1)^2$ has an obvious one-square certificate, which the
checker accepts when explicitly supplied. The restricted LP dictionary fails
to find a certificate for it. This failure is expected: a rank-one quadratic
Gram matrix for the coefficient vector $(1,1,1)$ cannot be decomposed as a
positive sum of rank-one vectors with only two nonzero coordinates, unless the
summands lie in its one-dimensional range; none of the allowed nonzero
binomial vectors do.

Thus the \emph{checker} supports a larger certificate language than this
\emph{searcher} can discover. The experiment is a useful protection against an
inflated ``SOS-complete'' claim. A production solver should escalate to general
Gram-matrix search or use the existing Lean SOS engine, rather than interpret
this failure as a negative mathematical result.

\subsection{Production search: reuse, rational reconstruction, and escalation}
The recommended nonlinear adapter first tries cheap sign reasoning and existing
arithmetic closers. It then tries a sparse dictionary search for a compact
certificate, followed by the compatible Lean SOS engine. Only afterward should
it increase degree, constraint-product, or refutation budgets. Those escalation
parameters already exist in the inspected SOS engine; Forge's contribution is
to allocate them according to outstanding proof obligations and reuse their
results.

For a general SOS block, write $\sigma=v^TQv$ in a monomial basis $v$, search
for a positive semidefinite Gram matrix $Q$, and recover rational data. A
weighted-square certificate can be extracted from an exact rational
factorization $Q=LDL^T$ with nonnegative diagonal $D$. Zero pivots require care:
a zero pivot is admissible only when the relevant residual off-diagonal entries
are zero, or after a justified permutation/reconstruction. A floating-point
``nearly positive semidefinite'' matrix is not adequate.

CSDP is a concrete open-source SDP search backend; the existing Lean SOS
integration is preferable to building a second unaudited wrapper.\footnote{The
COIN-OR CSDP project \cite{csdp}.}
A fallback oracle such as SymPy can suggest exact polynomial identities or
perform differential tests, but a returned identity must still be checked in
the chosen certificate representation.\footnote{SymPy polynomial tools
reference \cite{sympy}.}

To prove strict positivity, search for $p-\varepsilon\ge0$ with an exact
$\varepsilon>0$, or use a properly justified strict-constraint refutation.
A fixed global margin is sufficient but not necessary: $x>0$ implies $x>0$
without a uniform rational lower bound. For infeasibility, a certificate for
$-1$ in a nonnegative cone plus the equality ideal gives a contradiction.
If $q>0$ is a strict hypothesis, a checked identity expressing $-q^m$ as a
nonnegative cone plus vanishing ideal terms, with $m\ge1$, also gives a
contradiction. All strictness and exponent side conditions must be retained.

\subsection{Making nonlinear results useful to other engines}
A successful certificate should not be used only to close the current goal.
It can publish an inequality needed by a theorem application, a branch guard,
a witness obligation, or an induction invariant. Conversely, a linear solver
can publish proved bounds that tighten nonlinear search. The planner can rank
candidate nonlinear consequences by the number of waiting consumers and the
estimated cost of their certificates.

One concrete extension is to introduce an abstraction variable $z=xy$ together
with checked box bounds $\ell_x\le x\le u_x$ and
$\ell_y\le y\le u_y$. The four elementary product inequalities yield
\begin{align*}
 z&\ge \ell_xy+\ell_yx-\ell_x\ell_y,\\
 z&\ge u_xy+u_yx-u_xu_y,\\
 z&\le u_xy+\ell_yx-u_x\ell_y,\\
 z&\le \ell_xy+u_yx-\ell_xu_y.
\end{align*}
For example, the first follows from $(x-\ell_x)(y-\ell_y)\ge0$.
These are proved linear consequences once the box assumptions and $z=xy$ are
established. A linear model violating $z=xy$ can guide refinement, but it is
not a counterexample to the original goal. This bound/cut exchange is proposed;
it is not implemented in the supplied prototype.

\section{Exact Bernstein certificates on bounded domains}
\label{sec:bernstein}

\subsection{The representation and bound}
For a box $B=\prod_i[a_i,b_i]$ with rational endpoints and $a_i<b_i$, substitute
$x_i=a_i+(b_i-a_i)t_i$. Express the resulting polynomial on $[0,1]^n$ as
\[
 p=\sum_{0\le k_i\le d_i}\beta_k
       \prod_{i=1}^n\binom{d_i}{k_i}t_i^{k_i}(1-t_i)^{d_i-k_i}.
\]
Every basis function is nonnegative and the basis functions sum to one.
Consequently, $\min_k\beta_k\le p\le\max_k\beta_k$ on the box.

If the power-basis coefficients after substitution are $a_j$, the implemented
conversion is
\[
 \beta_k=\sum_{j\le k} a_j
          \prod_i\frac{\binom{k_i}{j_i}}{\binom{d_i}{j_i}}.
\]
All arithmetic uses exact rationals. A leaf certificate stores a claimed lower
bound, but the checker recomputes it; the claim itself is never trusted.

\subsection{A covering tree, not a collection of sampled boxes}
When a leaf bound is insufficient, bisect its longest side at a rational
midpoint. An internal certificate node records the split axis and split value
and has exactly two children. The checker reconstructs the two child boxes
from the parent, verifies that the split is interior, and checks both children.
This inductive construction proves coverage automatically. Accepting a list
of independently positive boxes without checking that they cover the original
domain would leave a soundness gap.

\begin{proposition}[Box-certificate soundness]
If every leaf of a checked covering tree has a nonnegative Bernstein lower
bound, then $p\ge0$ throughout the original box. If every leaf bound is strictly
positive, then $p>0$ there.
\end{proposition}
\begin{proof}
The Bernstein convex-combination identity proves the leaf claims. Each internal
node's children cover its box, including their common boundary. Induction over
the tree transports the property to the root.
\end{proof}

\subsection{An executed example and a termination qualification}
For $p(x)=x^2-x+3/10$ on $[0,1]$, the initial Bernstein coefficients are
\[
 (3/10,-1/5,3/10).
\]
Their negative minimum is inconclusive, not a refutation. Splitting at $1/2$
produces two leaves, each with exact lower bound $1/20>0$. The certificate has
three nodes and proves strict positivity on the requested interval.

The polynomial is actually positive on all of $\R$; an SOS proof is simpler.
The bounded-domain experiment is intentionally used to exercise a different
certificate mechanism. It should not be advertised as a case that demands
Bernstein subdivision or defeats the SOS engine.

For a polynomial with a positive minimum on a compact box, sufficiently fine
uniform subdivisions yield positive Bernstein lower bounds. To see the
mechanism, expand around a subbox endpoint. Under affine rescaling, every
nonconstant coefficient contains at least one side-length factor; the
correction of a Bernstein coefficient from the endpoint value tends uniformly
to zero as all side lengths tend to zero. A positive minimum eventually
absorbs that correction. This explains eventual success with unbounded fair
subdivision, not any useful a priori complexity guarantee.

If the minimum is zero, the same argument gives no finite strict certificate;
a true non-strict inequality can also remain hard for a chosen subdivision
policy. Production Forge should try symbolic factorization, exact boundary
analysis, or another certificate language before subdividing indefinitely.
The tensor-product basis has $\prod_i(d_i+1)$ coefficients, so dimension is
an important practical limit.

\section{Theory boundaries, SAT, and higher-order reasoning}
\label{sec:boundaries}

\subsection{Typed bridges rather than erased semantics}
Sharing proved propositions is safer than assuming every solver's internal
objects have interchangeable meanings. Each reifier should return a typed
interpretation, an atom environment, and a proof that the encoded target denotes
the original expression. If two solvers use different abstractions of a term,
exchange a checked bridge theorem or retain the representations separately.

For example, rational polynomial algebra treats subtraction as additive
inversion. Natural-number subtraction is truncated. Translating
\code{(a - b : Nat)} to a field difference requires a guard such as $b\le a$
and an appropriate transport theorem; otherwise one must keep the operation
opaque or split the relevant cases. Integer division is not field division,
and fixed-width addition is not unbounded integer addition. These distinctions
should be visible in the engine capability descriptors.

Unknown real-valued applications can be polynomial atoms without being
uninterpreted everywhere. One engine may know a theorem about $f(x)$ while the
polynomial engine treats its value as one coordinate. Once that theorem gives
a bound or equality, the bound is proved in the original context and can be
translated into the polynomial representation. No global erasure of dependent
types is necessary.

\subsection{A conservative higher-order adapter}
Extensionality is already present in the baseline. Forge's higher-order plan
is to choose useful extensionality steps, expose resulting pointwise demands,
and use a reconstruction-capable Duper/Lean-auto route when appropriate.
Bound-variable identity, universe levels, and typeclass instances must remain
part of theorem matching. Matching under binders should use Lean's own typed
expression machinery, not the Python prototype's first-order matcher.

Monomorphization and translation to first-order or higher-order external
languages need explicit fragment checks and proof-producing transport. Empty
types, dependent function domains, and quantifier well-formedness must not be
silently replaced by the external solver's assumed nonempty sorts. Unsupported
constructs should remain opaque where a sound abstraction permits that, or
cause the adapter to abstain.

\subsection{SAT and conflict learning}
For Boolean or fixed-width bitvector obligations, a SAT-oriented backend may be
much more suitable than broad quantified saturation. Forge can also use a
checked conflict clause as a learned fact. If a branch refutation depends on
assumptions $A_{i_1},\ldots,A_{i_r}$, derive the corresponding clause
$\neg A_{i_1}\lor\cdots\lor\neg A_{i_r}$ with its proof. Conflict analysis may
suggest a smaller support set, but replay must justify the actual clause.

A proposed future CDCL-style coordinator can use these clauses to avoid
repeated branches. The first implementation should reuse existing solver
reasoning and branch-state machinery instead of writing another full SAT/SMT
solver. The included prototype has no CDCL implementation, bit-blaster, or
higher-order prover.

Do not assume a generic Nelson--Oppen completeness theorem simply because
solvers exchange equalities. Finite datatypes and fixed-width bitvectors do not
satisfy the usual stably-infinite assumptions, and dependent encodings add
further issues. The proof-producing architecture supports sound exchanges
without claiming a blanket completeness result for the combined theory.

\section{Trust profiles and final acceptance}
\label{sec:trust}

\subsection{Kernel checking and native evaluation are different profiles}
The current Lean tactic reference explicitly distinguishes kernel reduction
from native evaluation: \code{decide +native} and \code{native\_decide} enlarge
the trusted base, and the \code{bv\_decide} documentation notes code-generator
trust and generated axioms. Merely caching an LRAT file for \code{bv\_check}
does not, by itself, establish a kernel-only replay profile.\footnote{Lean
language reference, decision procedures and SAT integration \cite{tactics}.}

Forge should therefore have two explicit profiles. The default
\emph{standard-axiom, kernel-replay profile} permits the project's declared
logical axioms, but no newly admitted solver or native-evaluation assertions.
An optional \emph{native-evaluation profile} can accept a broader, explicitly
reported dependency policy. These labels concern logical trust, not whether
untrusted search is implemented in native code: an external optimizer is
perfectly compatible with strict replay when its returned data is checked
without trusting the optimizer.

For a certified computation, prove the checker's soundness theorem once, then
establish the concrete check with a supported kernel reduction path, such as
\code{decide +kernel}, or with an ordinary proof term. The SOS source's
proof-facing design is useful here. For a SAT checker, an adapter still has to
verify that its selected replay route meets the requested profile; it should
not claim a strict route exists merely because a tactic has ``check'' in its
name.

\subsection{An explicit acceptance procedure}
An eventual \code{forge} invocation should perform the following final checks:
\begin{enumerate}
\item Restore the recorded original goal context and instantiate the selected
proof's metavariables. Reject any unresolved term or universe metavariables.
\item Infer and check the proof's type against the exact original target,
using only legitimate definitional equality or recorded proof transports.
\item Inspect the transitive declaration dependencies. Reject
\code{sorryAx}, temporary admitted hypotheses, unapproved native assertions,
and any axiom outside the declared policy. Existing project axioms must be
reported, not quietly relabeled as part of the kernel.
\item Validate the assembled declaration through Lean's normal checking path.
For reproducibility, replay the exported proof in a clean pinned workspace
with the original pre-theorem environment.
\end{enumerate}

Type inference in a running metaprogram is useful but is not the whole audit.
An expression may type-check because an imported axiom states the desired
result. Likewise, a tactic can appear to close the metavariable while its
suggested source script leaves a hint unproved. Both require scrutiny beyond
the visible number of remaining tactic goals.

\subsection{Python certificates are not Lean proofs}
The Python checker suite shares low-level term and rational-polynomial code
with the search programs, though its replay algorithms do not call the search
procedures. Differential tests against SymPy reduce the risk of some arithmetic
implementation bugs, but they do not formally verify Python, its runtime, or
the data model. The correct claim is \emph{independently replayed exact
certificates for a specified fragment}, not kernel-checked Lean theorems.

The command \code{python -S scripts/verify\_certificates.py} is useful evidence
of separation: it replays persisted certificates without importing SciPy,
NumPy, or SymPy. It is not a replacement for porting the checker and its
soundness proof into Lean. The JSON parser also is not a hardened untrusted
network service; production deployments need size, nesting, time, and memory
limits in addition to mathematical validation.

\section{Scheduling, retrieval, and proof-size control}
\label{sec:scheduling}

\subsection{A deterministic initial schedule}
Start with a deterministic policy that can be reproduced and ablated. Give
cheap normalization and the pinned baseline a short initial opportunity.
Thereafter maintain separate queues for local saturation, theorem demands,
structural plans, witness candidates, nonlinear certificates, and external
logical search. Use fixed quanta plus a small fair-widening quota. Exact numeric
weights should be tuning parameters, not hard-coded claims about all goals.

One reasonable initial policy ranks actions by
\[
 \frac{\text{estimated probability of closing an obligation}
       \times(1+\text{waiting consumers})}
      {\text{estimated search cost}+\text{estimated replay cost}+1},
\]
subject to periodic service for each enabled queue. Initial estimates can be
simple tables by goal class; no trained model is required. A learned ranking
model may later improve efficiency, but it remains an untrusted heuristic and
cannot modify the proof acceptance rules.

Record a budget vector, not just a timeout: Lean heartbeats, external wall time,
maximum term generation, theorem instances, induction depth, candidate grammar
size, polynomial degree, certificate term count, coefficient bit length, and
replay steps. Every expensive engine should return an explicit exhaustion reason.
A mathematical failure and a resource failure must not have the same diagnostic.

\subsection{Library retrieval driven by missing premises}
The baseline already has a suggestion mechanism. Forge should improve how
retrieval is used, not claim to introduce library search. Query a theorem index
with both the goal and currently blocked premises. Prefer type-compatible
conclusion heads, operations appearing in the residual, and theorems whose
premises have likely producers in the active engines.

Use staged top-$k$ retrieval: a small set first, then larger sets if search
stalls. In a strict benchmark, build the index only from the original
pre-theorem environment. Never inspect the target theorem's existing proof to
choose premises unless the experiment is explicitly labeled an oracle-premise
upper bound. Attributes that indirectly reintroduce a held-out theorem are
also leakage.

Retrieval can propose an induction lemma, a bridge theorem, or a normalization
identity. It need not add every retrieved theorem to every solver. Record a
small premise slice for each engine call, then minimize that slice after a
successful proof if doing so makes replay materially cheaper.

\subsection{Avoiding repeated work and invalid caches}
A cache key must include the environment identity, original typed expressions,
local-context identity, universe instantiations, transparency policy where
relevant, and the dependency support of branch assumptions. Alpha-normalized
pretty-printed goals are not enough. A proof of $P(x)$ under $h:x=0$ is not a
proof of the same displayed goal under $h:x=1$.

For an initial implementation, restrict cross-branch sharing to closed theorems
or explicitly abstracted proofs. Share more aggressively only after context
transport is implemented and tested. Persist certificate data and source-level
replay recipes, not arbitrary serialized mutable metavariable contexts.

Proof construction also has a budget. Hash-cons repeated certificate subterms,
share derived lemmas, and prefer a compact algebraic identity over thousands of
small rewrites when its checking cost is lower. Preserve a readable proof
script where practical, but retain the exact certificate when a short script
would rerun nondeterministic search.

\section{Concrete Lean implementation plan}
\label{sec:implementation}

\subsection{Integration points that actually exist}
The pinned source contains \code{Lean.Meta.Grind.main}, parameter constructors
including \code{mkDefaultParams}, and a result structure with failure and
instrumentation information. There is an extra-fact path that consumes proofs.
These give a concrete starting point for an adapter.\footnote{Lean 4.34.0,
\code{Lean/Meta/Tactic/Grind/Main.lean} \cite{grindmain}.}

The internal state includes proof-bearing new facts, equality information,
generations, and saved metaprogramming state. It is substantial reusable
infrastructure, but it is not evidence of a stable public plugin interface for
all of Forge's proposed operations.\footnote{Lean 4.34.0,
\code{Lean/Meta/Tactic/Grind/Types.lean} \cite{grindtypes}.}

The safest first version operates at coarse granularity: construct local Lean
subgoals, run selected engines on those goals, collect closed proof terms, and
restart or re-enter \code{grind} with the proved consequences. This costs more
than a fully incremental integration but avoids unsafely mutating undocumented
internal structures. A later version can add narrow upstream hooks for
certified fact insertion, demand notification, resumption, and diagnostics.

\subsection{Proposed module layout}
\begin{longtable}{p{0.27\linewidth}p{0.65\linewidth}}
\toprule
Module & Responsibility and first implementation choice \\
\midrule
\endhead
\code{Forge.Core} & Session identities, typed facts, obligations, dependencies, budgets, and event vocabulary. \\
\code{Forge.Check} & Exact-target checking, axiom policy, certificate replay registration, and final export. \\
\code{Forge.Search} & AND/OR scheduler and transactional branch snapshots; borrow Aesop design patterns or integrate suitable internals behind an adapter. \\
\code{Forge.Demand} & Conclusion index, typed backward matching, premise joins, and demand/forward coordination. \\
\code{Forge.Induction} & Recursive-position analysis, dependency-aware generalization, recursor application, and residual extraction. \\
\code{Forge.Lemma} & Typed grammar enumeration, safe test filters, anti-unification proposals, and the checked theorem bank. \\
\code{Forge.Witness} & Typed candidate terms and generation of exact witness obligations. \\
\code{Forge.Nonlinear} & Adapters for cheap arithmetic, compatible SOS engine, and the polynomial/box certificate routes. \\
\code{Forge.Bridges} & Proved representation conversions and their guards; no global semantic erasure. \\
\code{Forge.Backends} & Versioned Grind, Duper, Lean-SMT/QuerySMT, and optional SAT interfaces. \\
\code{Forge.Replay} & Stable proof scripts/certificate files, dependency manifests, and deterministic replay without numerical search. \\
\bottomrule
\end{longtable}

The separation between \code{Core}, \code{Check}, and external-search modules
allows replay-only builds. A numerical solver can be absent on the machine
checking a finished proof. Compatibility wrappers belong in their own modules;
a Lean upgrade should not require editing every search algorithm.

\subsection{A proposed user interface}
The following syntax expresses intended behavior only. It is not accepted by
the supplied files, which do not define a \code{forge} elaborator.
\begin{Code}
by
  forge

by
  forge (config := {
    induction := true
    synthesizeLemmas := true
    nonlinear := true
    externalSearch := false
    trust := .standardAxioms
  }) [helper_theorem]

by
  forge?  -- suggest a replayable script/certificate invocation
\end{Code}

A failure report should name the smallest remaining obligations and the
exhausted budgets. For example: ``Induction step requires a lemma about the
changing accumulator; 128 candidate expressions rejected by tests; 3 universal
proofs remain open.'' Another useful report is: ``SOS candidate rejected by
exact coefficient check; no theorem added.'' Such messages help distinguish a
missing library theorem, a poor search bound, and a genuine mathematical issue.

\subsection{Implementation stages with acceptance gates}
\textbf{Stage A: trustworthy orchestration.} Implement the state model,
transactional goal isolation, coarse \code{grind} adapter, fact support,
exact-goal validation, and replay export. Acceptance requires regression tests
for sibling-branch leakage, metavariable contamination, incorrect transports,
and denied axioms. A successful \code{grind} proof should survive orchestration
without logical changes.

\textbf{Stage B: the first major capability extension.} Port the list-fragment
induction checker or generate ordinary Lean induction proofs, then implement
recursive-position motive selection and changing-parameter generalization.
Add the bounded lemma grammar. Acceptance requires solving unseen variations
of accumulator correctness with the correctness theorem absent from the
environment, together with false-statement and circular-lemma rejection tests.

\textbf{Stage C: theory and witness integration.} Align a specific SOS/Hex/Mathlib
manifest with the target Lean version; expose certified nonlinear consequences
as reusable facts; connect witness enumeration and guard demands. Add a Lean
Bernstein soundness theorem and exact covering-tree replay if benchmarks justify
it. Acceptance requires certificates replaying without a numerical solver and
mixed examples needing both structural and arithmetic reasoning.

\textbf{Stage D: wider quantified and external search.} Add the demand index,
semi-naive premise joins, and optional reconstruction-capable external engines.
Add a strict SAT path only when its checker and trust policy are established.
Acceptance requires zero unproved hints in all counted successes, explicit
compatibility manifests, and comparisons against the strongest configured
baseline rather than default \code{grind} alone.

\textbf{Stage E: optimize from measurements.} Introduce persistent incremental
hooks, learned scheduling, stronger lemma synthesis, and richer algebraic
certificates only after profiling demonstrates their value. Measure replay
cost as well as search cost. These are deliverable stages rather than promised
calendar estimates; the difficult work lies in dependent-context handling,
version alignment, and proof reconstruction, not in writing the tactic keyword.

\section{Executed experiments and reproducibility}
\label{sec:experiments}

\subsection{Environment and measured scope}
The recorded run uses Python 3.13.5 on Linux x86-64, NumPy 2.3.5, SciPy 1.17.0,
SymPy 1.14.0, and pytest 9.0.2. The exact replay modules use the standard library;
NumPy/SciPy are used by numerical search, and SymPy is used for differential
arithmetic tests. The saved environment and run output are in
\code{results/results.json} and \code{results/experiment\_log.txt}.

The suite completed with \textbf{69 passing tests}. Thirteen persisted
certificate bundles then replayed successfully in a separate command run with
Python's site initialization disabled. One bundle contains an ordered bank of
five inductively proved lemmas and a specialization proof; four contain Horn
proofs; six contain cone certificates; one contains a box certificate; and one
contains the two certificates for an existential witness. These are bundles,
not 13 independent Lean theorems or 69 benchmark problems.

\begin{table}[htbp]
\centering
\caption{Executed mechanism results. ``Checked'' means the Python fragment checker.}
\begin{tabularx}{\linewidth}{p{0.28\linewidth}X}
\toprule
Mechanism & Recorded result \\
\midrule
Accumulator lemma discovery & Direct induction fails in the prototype; candidate 41 yields the generalized specification; induction and specialization check. \\
Additional list obligations & Five total lemmas receive checked induction certificates in dependency order. \\
Quartic nonnegativity & 280 LP columns, 35 equations, three exact weighted-square terms. \\
Guarded nonnegativity & Four additional targets certified using their explicit guards. \\
Incomplete dictionary control & Search fails on $(x+y+1)^2$; a separately supplied one-square certificate checks. \\
Bernstein box reasoning & Initial bound inconclusive; one split yields two positive leaves and a three-node certificate. \\
Polynomial witness & Candidate 23 gives $x^2+1$; both universal inequalities check. \\
\bottomrule
\end{tabularx}
\end{table}

\begin{table}[htbp]
\centering
\caption{Horn policy comparison; neither column is an implementation of \texttt{grind}.}
\begin{tabular}{rrrrl}
\toprule
Depth & Forward facts & Backward proof nodes & Forward cap & Forward status \\
\midrule
4 & 16 & 5 & 5,000 & Certified \\
8 & 256 & 9 & 5,000 & Certified \\
12 & 4,096 & 13 & 5,000 & Certified \\
16 & 5,000 & 17 & 5,000 & Cap reached \\
\bottomrule
\end{tabular}
\end{table}

Raw single-run timings are retained for reproducibility, but no speedup ratio
is used as an article conclusion. The examples are tiny, numerical libraries
have initialization effects, and different search policies perform different
amounts of work. Node counts and checked certificates are the more interpretable
observations here.

\subsection{Negative tests are part of the evidence}
The tests reject negative cone weights, absent guard indices, wrong polynomial
dimensions, and a target changed by only $10^{-12}$. They reject an out-of-range
box split, an incorrect leaf bound, a missing child, and a strict claim based
on a zero lower bound. A false polynomial target produces no certificate.

Horn mutation tests reject cyclic parent indices, the wrong rule, the wrong
final goal, and a removed premise. A guarded rule cannot fire when its guard is
missing. The backward engine's unsupported unbound-premise case is tested to
abstain. Induction tests reject a base-case use of the induction hypothesis,
invalid generalization metadata, external rigid-variable intrusion, a false
equality, and a reference to an unavailable future lemma. Removing required
accumulator generalization breaks the certificate as expected.

Twelve seeded tests compare sparse arithmetic, evaluation, and substitution
against SymPy, and twelve seeded planted-cone tests check numerical search
followed by exact replay. These are useful robustness checks, not an exhaustive
proof of the implementation. The shared polynomial primitives remain part of
the Python experiment's trusted implementation base.

\subsection{Reproduction commands}
From the extracted archive root:
\begin{Code}
python -m pip install -r requirements.txt
python -m pytest -q
python scripts/run_experiments.py
python -S scripts/verify_certificates.py
python scripts/check_lean.py
\end{Code}
The last command records \code{NOT\_RUN} when Lean is absent; it must not be
interpreted as a successful Lean test. When a suitable Lean executable is
available, it attempts to compile the core structural reference file and records
the actual version and output. The separate Mathlib examples need a compatible
Mathlib project and are not compiled by that script.

The article can be rebuilt with:
\begin{Code}
cd article
xelatex -interaction=nonstopmode -halt-on-error forge.tex
xelatex -interaction=nonstopmode -halt-on-error forge.tex
xelatex -interaction=nonstopmode -halt-on-error forge.tex
\end{Code}
No external bibliography processor is necessary. Font programs are not included;
the document uses fonts available in common TeX/font installations. The source
and code are UTF-8.

\section{A fair end-to-end evaluation of ``significantly more capable''}
\label{sec:evaluation}

\subsection{Why the environment and completion metric matter}
The QuerySMT study explicitly discusses both a success metric that can leave
individual hints unproved and contamination from benchmark environments that
already contain target theorems accessible through attributes. Those results
are valuable but cannot be copied directly as fully verified current-version
success rates.\footnote{QuerySMT methodology and evaluation discussion,
Sections 8.1--8.2 \cite{querysmt}.}

Forge should count a solved problem only when the exported proof checks in the
pre-theorem environment and passes the declared axiom policy. Adding a theorem
alias in a file that imports the original theorem is not a sufficient
holdout strategy. The test harness should construct the declaration environment
before the target is introduced, or reproduce a source prefix with an
appropriately isolated tactic import layer.

\subsection{Benchmark families and baselines}
Use separate tracks for ordinary mathematical developments and deliberately
minimal libraries. The first measures practical utility with normal imports;
the second measures structural and lemma-synthesis capability when the exact
answer is absent. Include recursive list/tree/accumulator correctness,
quantified equality and Horn chains, nonlinear inequalities, existential
witnesses, and mixed structural/arithmetic obligations. Preserve false and
ill-typed near-neighbors as soundness regressions, not as ordinary success-rate
items.

Compare at least: pinned default \code{grind}; \code{grind} with suggestions
and reasonably tuned budgets; a simple Aesop/Grind portfolio; compatible Duper
and reconstruction-based SMT routes; and a basic portfolio with the existing
SOS tactic. Otherwise a gain attributable merely to adding a known solver may
be mistaken for a gain from Forge's coordination.

Use the same pre-theorem declarations, permissible retrieved premises, total
wall-clock budget, memory limit, and trust profile. External solver time,
translation, reconstruction, proof export, and replay must all be charged.
Report cold-start and warm-cache conditions separately. Tune on development
theorems; hold out entire modules or definition families for testing so near-duplicate
lemmas do not dominate the results.

\subsection{Metrics, ablations, and a concrete acceptance target}
Report solved counts by family, uniquely solved cases, failures by reason,
median and tail latency, proof bytes, certificate bytes, reconstruction time,
and checking time. A conditional proof with one open guard belongs in a
partial-progress category. A model of an abstraction belongs in a diagnostic
category until it is validated against the original statement.

Ablate induction planning, lemma synthesis, demand propagation, witness
synthesis, nonlinear fact exchange, and proof-cost-aware scheduling separately.
Most importantly, compare the full planner against the \emph{same engines run
without feedback}. That identifies whether the architecture contributes more
than a portfolio union.

A useful engineering gate would require a predeclared improvement on at least
two structurally different held-out families, together with a bounded regression
rate on easy baseline successes and no trust-policy violations. The numerical
thresholds should be set before benchmarking, for example a target of 20\%
more fully checked solves on the selected difficult families and at most 5\%
regression on the baseline suite. These numbers are proposed acceptance
criteria, \emph{not achieved results or forecasts}. No such Lean evaluation has
been run for this artifact.

\section{Limits, priorities, and conclusion}
\label{sec:conclusion}
The largest practical risk is search explosion. Induction choices multiply
subgoals; lemma grammars grow exponentially; higher-degree polynomial
certificates become large; backward quantifier search can branch as badly as
forward saturation. A planner makes these choices explicit and budgeted, but
does not remove their underlying difficulty. Completeness claims must always
refer to a specified fragment and a specified unbounded search policy.

The second risk is integration rather than mathematics: mutable metavariable
state, dependency transport, representation semantics, version mismatch, and
large proof terms. The staged design prioritizes these boundaries before
advanced heuristics. A restricted correct implementation that exports a small
set of ordinary Lean proofs is preferable to a wide solver bridge that silently
leaves obligations or enlarges the trust base.

The most valuable first release is therefore not ``every solver behind one
keyword.'' It is a planner that automatically discovers and proves useful
induction strengthenings, synthesizes witnesses whose obligations can be
certified, and coordinates those results with the already strong Grind, SOS,
and logical engines. The Python experiments show that these mechanisms can be
made concrete, inspectable, and certificate-driven. They do not establish that
a finished Forge tactic outperforms current Lean automation.

\textbf{Bottom line.} There is a technically credible route to substantially
broader automation than a local \code{grind} call: let the tactic search for the
missing proof structure, not only for more consequences of a fixed local
state. The next decisive milestone is a Lean implementation of the scoped
planner and structural synthesis path, followed by clean-environment,
fully-checked comparisons against a strong multi-engine baseline.

\appendix
\section{Artifact map and status ledger}
\begin{longtable}{p{0.41\linewidth}p{0.51\linewidth}}
\toprule
Path & Contents and validation status \\
\midrule
\endhead
\code{article/forge.tex}, \code{forge.pdf} & This article and its rendered PDF. \\
\code{forge\_lab/terms.py} & Typed first-order terms, fixed list equations, rewrite search and independent replay. Executed. \\
\code{forge\_lab/induction.py} & Generalization-specific lemma enumeration and list induction certificates. Executed. \\
\code{forge\_lab/horn.py} & Two bounded Horn search policies and a rule-instance DAG checker. Executed. \\
\code{forge\_lab/poly.py} & Exact sparse rational arithmetic, cone checker, Bernstein transformation, and covering-tree checker. Executed. \\
\code{forge\_lab/nonlinear.py} & Untrusted SciPy LP search and bounded exact box search. Executed. \\
\code{forge\_lab/witness.py} & Polynomial witness enumeration with two universal certificates. Executed. \\
\code{tests/test\_certificates.py} & 69 tests, including arithmetic differential tests and deliberate certificate corruptions. Passed. \\
\code{results/} & Raw outputs, JSON certificates, count CSV, test/replay logs, and explicit Lean nonexecution status. \\
\code{scripts/} & Experiment, standard-library replay, and optional Lean compilation commands. \\
\code{lean/Structural.lean} & Ordinary core Lean proofs corresponding to the structural examples. Supplied, not compiled. \\
\code{lean/CertificateExamples.lean} & Mathlib proof scripts illustrating the arithmetic certificates. Supplied, not compiled. \\
\code{sources.json} & Bibliographic and source-version inventory with original URLs. \\
\bottomrule
\end{longtable}

\section{Certificate schemas and their interpretation}
A cone record contains the exact target polynomial, an ordered list of guards,
and terms $(c,q,J)$. The interpretation is
$p=\sum c q^2\prod_{j\in J}g_j$, conditional on each listed guard being
nonnegative. All rational values are serialized as strings, never floating-point
JSON numbers. Variable count and monomial exponent vectors are explicit.

A box record contains the exact polynomial, the root box, a strictness flag,
and a tree of leaves or interior splits. The checker reconstructs child boxes;
there is no trusted user-supplied coverage assertion. A witness record contains
one univariate polynomial and cone certificates for its difference from $x$
and from $-x$.

A Horn record contains the original premises, named rules, exact target, and
proof DAG. Its validity is conditional on the listed premises and rules; those
inputs are the problem statement, not new discoveries. A list-induction record
contains the candidate equality, major variable, generalized variables, and
four rewrite traces: left/right traces for the base and step cases. The theorem
bank's order supplies the allowed previously checked lemmas.

The saved records are intended to be portable explanatory artifacts for the
prototype. They are not yet an agreed wire format for Lean. A production format
must additionally bind every object to a versioned typed interpretation and
context manifest, impose resource limits, and have a checked Lean decoder or
an equally well-justified elaboration path.

\section{A compact replay target for the discovered lemma}
The following intended Lean proof explains the desired final output of
structural synthesis. The full reference file defines the recursive operations
and preceding concatenation lemmas. This source was not compiled here.
\begin{Code}
theorem revAcc_spec (xs acc : List α) :
    revAcc xs acc = app (rev xs) acc := by
  induction xs generalizing acc with
  | nil => rfl
  | cons x xs ih =>
    simp only [revAcc, rev, ih, app_assoc, app]

theorem revAcc_correct (xs : List α) :
    revAcc xs [] = rev xs := by
  rw [revAcc_spec, app_right_nil]
\end{Code}
The final theorem does not invoke the synthesizer. It reuses a separately
proved general lemma. This is the preferred deployment pattern: potentially
expensive search once, then a stable and understandable proof for ordinary
recompilation.

\newpage
\renewcommand{\refname}{Sources}
\addcontentsline{toc}{section}{Sources}
\begin{thebibliography}{99}
\small
\bibitem{grind} Lean project. \emph{The grind tactic}. Lean Language Reference.
Observed September 14, 2026. \url{https://lean-lang.org/doc/reference/latest/The--grind--tactic/}.
\bibitem{release} Lean project. \emph{Lean 4.34.0 release}. September 14, 2026.
\url{https://github.com/leanprover/lean4/releases/tag/v4.34.0}.
\bibitem{config} Lean project. \emph{Grind configuration}, tag \code{v4.34.0}.
\url{https://raw.githubusercontent.com/leanprover/lean4/v4.34.0/src/Init/Grind/Config.lean}.
\bibitem{commring} Lean project. \emph{Grind commutative-ring constraint types}, tag \code{v4.34.0}.
\url{https://raw.githubusercontent.com/leanprover/lean4/v4.34.0/src/Lean/Meta/Tactic/Grind/Arith/CommRing/Types.lean}.
\bibitem{lia} Lean project. \emph{Linear Integer Arithmetic}. Lean Language Reference.
\url{https://lean-lang.org/doc/reference/latest/The--grind--tactic/Linear-Integer-Arithmetic/}.
\bibitem{aesop} Lean community. \emph{Aesop: White-box automation for Lean 4}. Repository documentation.
\url{https://github.com/leanprover-community/aesop}.
\bibitem{duper} Lean community. \emph{Duper}. Repository documentation.
\url{https://github.com/leanprover-community/duper}.
\bibitem{auto} Lean community. \emph{Lean-auto: Experiments on automation for Lean}. Repository documentation.
\url{https://github.com/leanprover-community/lean-auto}.
\bibitem{autopaper} Yicheng Qian, Joshua Clune, Clark Barrett, and Jeremy Avigad.
\emph{Lean-auto: An Interface between Lean 4 and Automated Theorem Provers}. 2025.
\url{https://arxiv.org/abs/2505.14929}.
\bibitem{smt} Lean-SMT project. \emph{Lean-SMT}. Repository documentation.
\url{https://github.com/ufmg-smite/lean-smt}.
\bibitem{querysmt} Joshua Clune, Haniel Barbosa, and Jeremy Avigad.
\emph{Hint-Based SMT Proof Reconstruction}. January 2026, arXiv:2601.14495v1.
\url{https://arxiv.org/html/2601.14495v1}.
\bibitem{sos} Lean SOS project. \emph{Harrison's sum-of-squares decision procedure for nonlinear real arithmetic in Lean 4}. Repository documentation.
\url{https://github.com/leanprover/sos}.
\bibitem{sosengine} Lean SOS project. \emph{SOS/Engine.lean}, observed \code{main} source, September 14, 2026.
\url{https://raw.githubusercontent.com/leanprover/sos/main/SOS/Engine.lean}.
\bibitem{soscore} Lean SOS project. \emph{SOS/Core.lean}, observed \code{main} source.
\url{https://raw.githubusercontent.com/leanprover/sos/main/SOS/Core.lean}.
\bibitem{sostoolchain} Lean SOS project. \emph{lean-toolchain}, retrieved \code{main} snapshot specifying \code{v4.32.0-rc1}; not a verified common build manifest.
\url{https://raw.githubusercontent.com/leanprover/sos/main/lean-toolchain}.
\bibitem{hex} Lean Hex project. \emph{hex-mv-poly: Mathlib-free executable multivariate polynomials for Lean 4}.
\url{https://github.com/leanprover/hex-mv-poly}.
\bibitem{linarith} Mathlib community. \emph{Mathlib.Tactic.Linarith.Frontend}. API documentation.
\url{https://leanprover-community.github.io/mathlib4_docs/Mathlib/Tactic/Linarith/Frontend.html}.
\bibitem{positivity} Mathlib community. \emph{Mathlib.Tactic.Positivity.Basic}. API documentation.
\url{https://leanprover-community.github.io/mathlib4_docs/Mathlib/Tactic/Positivity/Basic.html}.
\bibitem{leanegg} Marcus Rossel and contributors. \emph{lean-egg}. Repository, including current deprecation notice.
\url{https://github.com/marcusrossel/lean-egg}.
\bibitem{egg} E-graphs community. \emph{egg: e-graphs good}. Rust library repository.
\url{https://github.com/egraphs-good/egg}.
\bibitem{lemmasynthesis} Yican Sun, Ruyi Ji, Jian Fang, Xuanlin Jiang, Mingshuai Chen, and Yingfei Xiong.
\emph{Proving Functional Program Equivalence via Directed Lemma Synthesis}. 2024.
\url{https://arxiv.org/abs/2405.11535}.
\bibitem{highs} SciPy developers. \emph{linprog(method='highs')}. Optimization API documentation.
The executed artifact uses SciPy 1.17.0; the moving online documentation may differ.
\url{https://docs.scipy.org/doc/scipy/reference/optimize.linprog-highs.html}.
\bibitem{csdp} COIN-OR. \emph{CSDP: A solver for semidefinite programming problems}.
\url{https://github.com/coin-or/Csdp}.
\bibitem{sympy} SymPy developers. \emph{Polynomial Manipulation: Reference}.
\url{https://docs.sympy.org/latest/modules/polys/reference.html}.
\bibitem{tactics} Lean project. \emph{Tactic Reference}, especially decision procedures, native evaluation, and SAT integration.
\url{https://lean-lang.org/doc/reference/latest/Tactic-Proofs/Tactic-Reference/}.
\bibitem{grindmain} Lean project. \emph{Grind main entry points}, tag \code{v4.34.0}.
\url{https://raw.githubusercontent.com/leanprover/lean4/v4.34.0/src/Lean/Meta/Tactic/Grind/Main.lean}.
\bibitem{grindtypes} Lean project. \emph{Grind state and proof-bearing types}, tag \code{v4.34.0}.
\url{https://raw.githubusercontent.com/leanprover/lean4/v4.34.0/src/Lean/Meta/Tactic/Grind/Types.lean}.
\end{thebibliography}
\end{document}
